\chapter{AI Prompt Log}
\ChapterLabel{AIPromptLog}
\section{Chapter Responsibility}
\label{Section:AIPromptLogResponsibility}
Preserve the provenance of material AI assistance and its human review.

\noindent\textbf{Expected result:} An auditable account of how AI contributed to the engineering record.

\section{How to Use This Appendix}
Use this log when AI assistance materially affects requirements, architecture, code, tests, or release evidence. Record enough context for a reviewer to understand the input, output, assumptions, and human corrections.
\section{Initialization Interaction}
\label{Section:InitializationInteraction}
Date: November 12, 2026. Tool: Codex. Role: aiaCoordinator, with drafting and
self-review performed by the same agent. The user requested AIS-02 initialization
from an unchanged template copy, a facts/gaps/files/checks preview, confirmation
before dependent edits, separate TikZ sources, complete PDF build and inspection,
and a later acceptance/publication gate. The user then confirmed the project
name, namespace, owner and approver, documentation-only scope, Applicable
Embedded Systems profile, and user ownership of unknown technical details.
See \hyperref[Section:ConfirmedProjectDecisions]{confirmed decisions} and
\hyperref[Section:OwnedQuestions]{owned questions}. No product facts or evidence
were supplied. The implementation uses no delegated agents. At that stage,
owner acceptance, independent review, publication authorization, and repository
confirmation were pending. Subsequent owner disposition is recorded in
\hyperref[Section:PublicationAuthorization]{publication authorization}. Technical examples from the template were not adopted as project artifacts.

\section{Owner Identity Correction}
\label{Section:OwnerIdentityCorrection}
On November 12, 2026, \DocumentOwner\ supplied their name, organization
\Organization, and email \href{mailto:\DocumentEmail}{\DocumentEmail}, and
requested replacement of the generic owner fields. Codex updated the shared
metadata, front matter, PDF author, ownership references, and
\hyperref[Question:Q-01]{Q-01}. The change uses local branch
\texttt{codex/update-owner-details}; publication and its pull request remain
pending the owner's separate publication authorization at that time. Technical
questions and product release decisions remain open; subsequent project
acceptance is recorded below.

\section{Publication Authorization}
\label{Section:PublicationAuthorization}
On November 12, 2026, \DocumentOwner\ stated: ``I accept the initialized
project and authorize publication.'' The owner directed a source commit on
main and push to the repository identified in
\hyperref[Section:InitializationPublication]{the publication record}, with
branches and pull requests for every later change. Codex records this actual
human disposition and prepares the initial publication. It does not infer
technical approvals or product release from this authorization.

\section{Requirement Scope or Family Extension Record}
When an AI-assisted change proposes a requirement scope or family that is not
already defined, record the proposed uppercase code, full name, decision
boundary, rationale, representative examples, semantic-overlap review,
identifier-collision check, affected register and category-catalog changes, and
the authorized human's approval or rejection. Do not assign an identifier from
the proposed scope or family before approval.
